% ===== §10 Empirical Dimensions =====
\section{Empirical Dimensions}
\label{p16:sec:empirical}

\noindent
The theory of relational luxury makes several empirically tractable claims: that coupling degree, rather than economic value, predicts the relational consequences of gift-giving; that high-coupling-degree gifts generate GRW more effectively than low-coupling-degree gifts of equivalent or greater economic value; and that the recognition of coupling degree depends on a cultivable sensibility. This section proposes a research programme for testing these claims, while acknowledging---as the parallel section of \pinkref{Paper~XV} acknowledged \parencite{paperxv}---that the core phenomenon (coupling degree) is, like GRW itself, not directly measurable and must be approached through its conditions and consequences.

\subsection{Operationalising Coupling Degree}
\label{p16:subsec:operationalising}

\noindent
The coupling degree of an object with respect to a relational subject's spatiotemporal structure is not directly measurable, but its four dimensions can be operationalised through a combination of behavioural, self-report, and process measures.

\emb{Temporal depth} can be assessed through the documentation of the temporal investment embedded in the gift: the time elapsed between the decision to give and the giving (for sustained-investment gifts), or the immediacy of the giving relative to the moment of acquisition (for immediate-presence gifts). These temporal structures can be documented through retrospective interview and, where possible, through contemporaneous records.

\emb{Intentional intensity} can be assessed through the degree of specific attention to the recipient that the gift selection involved: whether the gift was chosen with reference to the recipient's specific characteristics, preferences, and relational history, or selected from a generic category by economic criteria. This can be measured through structured interviews that probe the selection process, coded for the presence and depth of recipient-specific consideration.

\emb{Non-substitutability} can be assessed through the counterfactual question of whether another object could have served equally well: the high-non-substitutability gift is one for which no economic equivalent would have carried the same meaning, while the low-non-substitutability gift is one that could have been replaced by any object of equivalent economic value. This can be measured through both giver and recipient reports of the gift's irreplaceability.

\emb{Existential investment} is the most difficult dimension to operationalise, but it can be approached through measures of the giver's attentional and emotional engagement in the gift process: the degree to which the giving involved genuine presence and care rather than delegated or routine acquisition. Physiological and behavioural measures of engagement during the gift process, combined with self-report measures of the giver's experience, can provide indirect indices of existential investment.

\subsection{The Coupling Degree and Relational Outcome Study}
\label{p16:subsec:outcome_study}

\noindent
The central empirical prediction of the relational luxury framework is that coupling degree, rather than economic value, predicts the relational consequences of gift-giving. This prediction can be tested through a study that independently measures the coupling degree and the economic value of gifts exchanged in intimate relations, and examines their relative contributions to relational outcomes.

\emb{Design}: A longitudinal study of gift exchange in intimate relations, in which couples document the gifts they exchange over an extended period, providing both economic value information and the coupling degree measures described above. Relational outcomes (relational quality, physiological synchrony, subjective wellbeing) are measured at intervals throughout the study period.

\emb{Predicted results}: The relational luxury framework predicts that coupling degree will be a stronger predictor of relational outcomes than economic value; that high-coupling-degree gifts will be associated with improvements in relational quality regardless of their economic value; and that high-economic-value but low-coupling-degree gifts (the proxy-purchased diamond) will be associated with weaker relational outcomes than low-economic-value but high-coupling-degree gifts (the flower on the path). The framework further predicts an interaction between coupling degree and recipient sensibility: the relational benefit of high-coupling-degree gifts will be greater for recipients with higher cultivated sensibility, because the recognition of existence-attestation requires the sensibility to perceive it.

\subsection{The Misrecognition Study}
\label{p16:subsec:misrecognition_study}

\noindent
The analysis of subject-embedding misrecognition generates testable predictions about the conditions under which receivers accurately or inaccurately perceive the coupling degree of gifts.

The framework predicts that receivers will systematically over-attribute coupling degree to high-economic-value gifts---that the sign exchange value of an expensive gift will induce the perception of existence-attestation even when the actual coupling degree is low. This induced misrecognition can be tested experimentally by presenting receivers with gifts that vary independently in economic value and coupling degree, and measuring their perceptions of the giver's existential investment. The framework predicts that economic value will inflate perceived existential investment independently of actual coupling degree---a measurable manifestation of the receiver's induced misrecognition.

The framework further predicts that this misrecognition will be moderated by sensibility: receivers with higher cultivated sensibility will be less susceptible to the inflation of perceived coupling degree by economic value, because their sensibility allows them to perceive actual coupling degree more accurately. This prediction connects the empirical study of misrecognition to the theory of sensibility as a cultivable recognitional capacity.

\subsection{Methodological Challenges}
\label{p16:subsec:methodological_challenges}

\noindent
The empirical study of relational luxury faces methodological challenges that parallel those identified in \pinkref{Paper~XV} \parencite{paperxv} and that must be honestly acknowledged.

\emb{The measurement problem}: Coupling degree, like GRW, is not directly measurable; all measures are indirect, inferring the coupling from its dimensions and its consequences. This indirection introduces systematic uncertainty into the empirical programme, and the interpretation of results must maintain the distinction between coupling degree and its measurable correlates.

\emb{The reactivity problem}: The measurement of coupling degree may itself alter the phenomenon being measured. Asking couples to document and reflect on the coupling degree of their gifts may increase their attentiveness to coupling degree, thereby changing their gift-giving behaviour and its relational consequences. This reactivity is particularly significant for the study of a phenomenon that depends, as relational luxury does, on attentive presence and reflective awareness.

\emb{The cultural specificity problem}: The operationalisation of coupling degree developed here is shaped by specific cultural assumptions about gift-giving, intimacy, and luxury. The cross-cultural application of these measures requires careful attention to the different cultural forms that coupling degree takes---the jade nourished by its wearer, the tea ceremony's \emph{ichi-go ichi-e}---and the development of culturally appropriate measures rather than the imposition of a single operationalisation across cultural contexts.

\emb{The self-report limitation}: Many of the measures of coupling degree rely on self-report, which is subject to the biases of social desirability, retrospective reconstruction, and the very misrecognition that the framework identifies. The giver who has misrecognised their own existence-attestation---who believes that their economic expenditure constitutes existential investment---will report high existential investment even when the actual coupling degree is low. The empirical programme must therefore triangulate self-report with behavioural and physiological measures that are less susceptible to this bias.

Despite these challenges, the empirical programme is worth pursuing, because the relational luxury framework makes predictions that diverge sharply from the predictions of the dominant economic-value framework, and the test of these divergent predictions would provide significant evidence about the nature of luxury in intimate relations. The prediction that the flower can outperform the diamond---that low economic value combined with high coupling degree can produce better relational outcomes than high economic value combined with low coupling degree---is, if confirmed, a result of considerable theoretical and practical significance.
